% =====================================================================
%  2bat-ud5-2.tex  —  SESSIÓ 2: Repàs amb dades reals
%  (sense preàmbul: s'inclou des de main.tex)
% =====================================================================
\horatitol{Sessió 2 — Repàs amb dades reals}%
{Activitat 1: sobre una taula de contingència real, calcular probabilitats simples i condicionades i interpretar-ne els patrons sense estigmatitzar.}

\subsection{Idea clau}

Recuperem tota la probabilitat de 1r en un cas real. Treballem sobre una \textbf{taula de
contingència}: una taula de doble entrada que creua dues característiques d'un col·lectiu.
D'aquí en traurem probabilitats \textbf{simples}, \textbf{compostes} i, sobretot,
\textbf{condicionades}, i n'interpretarem els patrons per orientar una intervenció.

\begin{idea}
\textbf{Llegir una taula de contingència.}
\begin{itemize}[leftmargin=1.4em, itemsep=3pt]
  \item \textbf{Probabilitat simple} $P(A)$: casos amb la característica $A$ dividit pel
        \textbf{total} de persones.
  \item \textbf{Probabilitat composta} $P(A\text{ i }B)$: casos que compleixen les
        \textbf{dues} característiques, dividit pel total.
  \item \textbf{Probabilitat condicionada} $P(A\,|\,B)$ (probabilitat de $A$
        \emph{sabent} $B$): es canvia la base — es divideix pel total \textbf{del grup
        $B$}, no pel total general.
\end{itemize}
\textbf{El pas delicat} és la condicionada: cal dividir pel total de la fila o columna que
fa de condició, no pel total de la taula.
\end{idea}

\noindent\textbf{Les dades.} Un servei ha atès $200$ persones. La taula creua el
\textbf{tipus de llar} amb si han \textbf{trobat feina} en un any:
\begin{center}
\renewcommand{\arraystretch}{1.3}
\begin{tabular}{l c c c}
\toprule
 & \textbf{Feina: sí} & \textbf{Feina: no} & \textbf{Total} \\
\midrule
\textbf{Viu sola}     & 45  & 35 & 80 \\
\textbf{Amb família}  & 90  & 30 & 120 \\
\midrule
\textbf{Total}        & 135 & 65 & 200 \\
\bottomrule
\end{tabular}
\end{center}

\subsection{Exemples}

\begin{exemple}[probabilitat simple i composta]
La probabilitat que una persona triada a l'atzar \textbf{visqui sola} és
\[
P(\text{sola})=\frac{80}{200}=\frac{2}{5}=0,4=40\%.
\]
La probabilitat que \textbf{visqui sola i hagi trobat feina} (les dues coses) és
$P(\text{sola i feina})=\dfrac{45}{200}=0,225=22,5\%$ (la casella d'encreuament, dividida
pel total).
\end{exemple}

\begin{exemple}[probabilitat condicionada: canviar la base]
Quina és la probabilitat que una persona hagi trobat feina \textbf{sabent que viu sola}?
Ara la base ja no és $200$, sinó les $80$ persones que viuen soles:
\[
P(\text{feina}\,|\,\text{sola})=\frac{45}{80}=\frac{9}{16}=0,5625\approx 56,3\%.
\]
En canvi, entre les que viuen amb família,
$P(\text{feina}\,|\,\text{família})=\dfrac{90}{120}=\dfrac{3}{4}=0,75=75\%$. La
diferència ($56\%$ enfront de $75\%$) és un \textbf{patró}: qui viu sol té, en aquest
col·lectiu, menys inserció laboral.
\end{exemple}

\subsection{Exercicis resolts}

\begin{exresolt}[condicionada amb la base canviada]
Calcula la probabilitat que una persona \textbf{visqui sola sabent que ha trobat feina}.
\solucio
La condició és «ha trobat feina»: la base són les $135$ persones amb feina. D'aquestes,
$45$ viuen soles:
\[
P(\text{sola}\,|\,\text{feina})=\frac{45}{135}=\frac{1}{3}\approx 0,333=33,3\%.
\]
Fixa't que $P(\text{sola}\,|\,\text{feina})=33,3\%$ \textbf{no} és igual a
$P(\text{feina}\,|\,\text{sola})=56,3\%$: condicionar en un sentit o en l'altre canvia la
base i el resultat.
\resposta{$P(\text{sola}\,|\,\text{feina})=\dfrac{1}{3}\approx 33,3\%$.}
\end{exresolt}

\begin{exresolt}[interpretar un patró sense estigmatitzar]
A partir de les dades, quin grup té més necessitat de suport per a la inserció laboral?
Com ho diries sense estigmatitzar?
\solucio
La inserció és més baixa entre qui viu sol ($56\%$) que entre qui viu amb família
($75\%$). Això indica que el grup que \textbf{viu sol} podria necessitar més
acompanyament en la cerca de feina. \textbf{Sense estigmatitzar:} no és que aquestes
persones «fallin», sinó que probablement disposen de menys xarxa de suport; la dada
serveix per \emph{orientar recursos} cap a qui més els necessita, no per etiquetar
ningú.
\resposta{El grup que viu sol; la lectura ha d'orientar recursos, no culpabilitzar.}
\end{exresolt}

\subsection{Exercicis per fer}

\noindent\textit{Fes servir la taula de contingència de dalt.}

\begin{exercicis}
\begin{exlist}
  \item Calcula $P(\text{amb família})$ (probabilitat simple) i expressa-la en fracció,
        decimal i percentatge.
  \item Calcula $P(\text{no feina})$: quina proporció del total no ha trobat feina?
  \item Calcula $P(\text{amb família i feina})$ (probabilitat composta).
  \item Calcula la condicionada $P(\text{no feina}\,|\,\text{sola})$: entre qui viu sol,
        quina proporció no ha trobat feina? \textbf{(Compte amb la base.)}
  \item \textbf{(Comparació)} Calcula $P(\text{no feina}\,|\,\text{sola})$ i
        $P(\text{no feina}\,|\,\text{família})$ i digues en quin grup és més alt el risc
        de no trobar feina.
  \item \textbf{(Interpretació)} A partir de les probabilitats condicionades que has
        calculat, redacta dues línies sobre cap a quin perfil convindria orientar el
        suport, tenint cura de no estigmatitzar.
\end{exlist}
\end{exercicis}

% ---------------------------------------------------------------------
\fullsolucions{Sessió 2}
\begin{solucions}
\begin{sollist}
  \item $P(\text{amb família})=\dfrac{120}{200}=\dfrac{3}{5}=0,6=60\%$.
  \item $P(\text{no feina})=\dfrac{65}{200}=\dfrac{13}{40}=0,325=32,5\%$.
  \item $P(\text{amb família i feina})=\dfrac{90}{200}=\dfrac{9}{20}=0,45=45\%$.
  \item Base: les $80$ que viuen soles; no feina $=35$.
        $P(\text{no feina}\,|\,\text{sola})=\dfrac{35}{80}=\dfrac{7}{16}=0,4375\approx 43,8\%$.
  \item $P(\text{no feina}\,|\,\text{sola})=\dfrac{35}{80}\approx 43,8\%$;
        $P(\text{no feina}\,|\,\text{família})=\dfrac{30}{120}=\dfrac{1}{4}=25\%$. El risc
        de no trobar feina és més alt entre qui viu sol ($43,8\%$ enfront de $25\%$).
  \item Resposta oberta. Idea: el suport a la inserció laboral es podria prioritzar cap a
        les persones que viuen soles, que mostren menys inserció i més risc; s'ha de
        formular com una orientació de recursos segons necessitat, no com un judici sobre
        les persones.
\end{sollist}
\end{solucions}
